\section{Висновки:}

В ході виконання лабораторної роботи було успішно досягнуто поставлену мету: 
розгорнуто кастомне робоче оточення для програмування на мові Java та опановано 
базовий синтаксис і оператори цієї мови через призму порівняльного аналізу з екосистемою C/C++.

На основі виконаних досліджень та розробленого програмного комплексу сформульовано такі науково-практичні висновки:
\begin{enumerate}
    \item \textbf{Ефективність інструментарію:} 
    Ручне налаштування універсального редактора VS Code та інтеграція компілятора JDK-27 
    за допомогою сценаріїв автоматизації Windows PowerShell (\texttt{tasks.json}) дозволили 
    створити прозорий та ефективний цикл збирання проєкту «з нуля». Такий підхід забезпечив 
    повний контроль над структурою пакетів та процесом компіляції без використання ваговитих 
    сторонніх систем збирання.
    \item \textbf{Концептуальні переваги Java:} 
    Практична реалізація алгоритмів підтвердила ключові архітектурні відмінності Java від 
    мов C/C++. Відсутність ручного керування пам'яттю та вказівників (завдяки роботі Garbage Collector), 
    строга ізоляція булевого типу та робота з масивами як з повноцінними об'єктами (властивість \texttt{length}) 
    суттєво підвищують безпеку розробки та нівелюють ризики виникнення критичних помилок переповнення буфера.
    \item \textbf{Об'єктно-орієнтована декомпозиція:} 
    На відміну від процедурного стилю, завдання лабораторної роботи були успішно розділені на ізольовані 
    класи-компоненти (\texttt{ArrayChecker}, \texttt{FizzBuzzGame}, \texttt{BalanceFinder}) із чіткою 
    інкапсуляцією логіки, що заклало правильний фундамент для подальшого вивчення об'єктно-орієнтованого програмування.
    \item \textbf{Культура розробки:} 
    Весь вихідний код програми був повністю адаптований до суворих вимог міжнародного стандарту \textit{Oracle Code Conventions}. 
    Впровадження трирівневої системи коментування (ліцензійний блок, JavaDoc-документація та внутрішні пояснення) разом із 
    автоматичною збіркою звіту в системі \texttt{LaTeX} дозволили вивести фінальну технічну документацію на рівень промислових стандартів розробки.
\end{enumerate}


\newpage
\section{Інформаційні ресурси та посилання на вихідний код}

Згідно з вимогами до здачі лабораторної роботи, повний проєкт вихідного коду розробленого 
додатку розміщено у віддаленому репозиторії системи контролю версій. 

\begin{itemize}
    \item \textbf{Посилання на Git-репозиторій проєкту:} \\
    \url{https://github.com/nasuehfw/oop-lab1.git}
    
    \item \textbf{Електронний каталог спільного доступу:} \\
    \url{https://drive.google.com/drive/folders/1mIulRQSfoHAyQ2GnqdQbwGz2dFX-5VGK?usp=drive\_link}
\end{itemize}
